% Related Work
% Organize by themes/categories, not paper-by-paper
% Show how your work builds on, differs from, or fills gaps in existing work

\label{sec:related}

% Introduction to related work
Research on Wikipedia spans multiple domains, including computer science, sociology, communication, and political science. Scholars have examined how Wikipedia governs itself, how collaborative editing impacts content quality, and how bias and neutrality emerge in open knowledge systems. This section organizes prior research into three major themes: governance and community culture, bias and neutrality, and democratic deliberation.

% Subsection 1: First category of related work
\subsection{Wikipedia Governance and Community Culture}

Reagle’s seminal work on Wikipedia culture~\cite{reagle2010good} explores the norms of collaboration and consensus that underpin the platform’s community ethos, often summarized by the term “good faith collaboration.” Konieczny~\cite{konieczny2009governance} and Shaw and Hill~\cite{shaw2014labor} expanded on this by identifying governance patterns that resemble oligarchic structures, where small, highly active groups control decision-making despite the open nature of the platform. Studies of editorial hierarchies show that informal authority and technical expertise influence whose edits persist and whose are reverted, shaping both policy enforcement and content stability.

\subsection{Bias and Neutrality in Wikipedia}
A second line of research examines whether Wikipedia’s open editing model fosters or mitigates bias. Hube~\cite{hube2017detecting} developed computational approaches to detect biased phrasing and framing in article text, showing that even factual statements can carry subtle ideological leanings. Greenstein and Zhu~\cite{greenstein2012ideological} compared political articles and found evidence of ideological balance emerging over time due to countervailing edits by users with opposing viewpoints. These studies underscore that bias in Wikipedia is dynamic—it evolves through ongoing revision and editorial negotiation.

\subsection{Democratic Deliberation and Online Collaboration}
Talk pages serve as deliberative spaces where editors debate sources, tone, and factual claims. Klemp and Forcehimes~\cite{klemp2010deliberation} highlight how these discussions mirror the ideals of deliberative democracy: open participation, reasoned argument, and consensus-seeking. Shi et al.~\cite{shi2019value} empirically demonstrated that politically diverse participation improves article quality and reduces ideological slant. However, participation remains uneven, and contentious topics often reveal asymmetrical power dynamics where a few editors dominate debates.

% Positioning your work
\subsection{Our Work in Context}
By combining machine learning–based bias detection with sentiment and toxicity analysis of Talk pages, this work investigates whether inclusive, democratic participation correlates with greater neutrality in Wikipedia’s content. This integrated perspective helps fill a key gap between computational text analysis and sociotechnical studies of online governance.
